\begin{document}
\title{\begin{tabular}[c]{@{}c@{}}\textbf{OOD Detection by an Orientation Likelihood Ratio}\end{tabular}}
\affiliate{sfc}{慶應義塾大学 総合政策学部B3 中澤・大越研究室 d-hacks}
\affiliate{sfc2}{慶應義塾大学大学院 政策・メディア研究科}
\author{omote}{}{sfc}[]
\author{親:zackeyさん}{}{sfc2}[]
\maketitle

\section*{概要}
凍結した拡散モデルによる分布外検知では, 1回の順伝播で中間特徴を読むCFSが最良であるが, 1{,}024次元の特徴を学習画像5万枚と比べる. 本研究の\Ours{}は, U-Netのブロックごとに数値を1つ(振幅)だけ読み,その並びを1つの式で判定する. 分布内の画像\olNRefMax{}枚の並びの平均と共分散,同じ画像を残り7つの向き(反転と回転)にしたときの並びの平均を求める. スコアは,分布内の平均を,向きを変えた平均の反対側へ動かした点へのマハラノビス距離で,共分散の比を固定した2つの正規分布の対数尤度比に等しい. 勾配による学習,追加のネットワーク,人手のラベル,分布外の画像は使わない. 5つの分布内データセットで,自然画像の分布外集合に対する平均AUROCは\olTestOL\%で,先行研究の最良(CFS)を\olGapCfsKnn{}ポイント上回った(95\%信頼区間[\olGapCfsKnnLo, \olGapCfsKnnHi]). 6つのテスト分割では\olSixMin{}〜\olSixMax\%であった. 順伝播は途中で止める1回で, CPU時間はCFSの\olSecRatioCfsKnn{}倍である. 縮小の履歴をそろえた対照でも差は変わらない. 定数1つは開発用の分割で決め,テスト分割は以前に他の手法で結果を見たものである.
